\documentclass[11pt]{article}
\usepackage[spanish]{babel}
\usepackage[utf8]{inputenc}
\usepackage[T1]{fontenc}
\usepackage{geometry}
\geometry{margin=2.5cm}
\usepackage{amsmath, amssymb, amsthm}
\usepackage{enumitem}
\usepackage{booktabs}
\usepackage{natbib}
\usepackage{hyperref}
\hypersetup{colorlinks=true, linkcolor=blue, citecolor=blue, urlcolor=blue}

\theoremstyle{definition}
\newtheorem{supuesto}{Supuesto}
\newtheorem{hipotesis}{Hipótesis}
\theoremstyle{plain}
\newtheorem{prop}{Proposición}

\title{Marco teórico tentativo\\
\large Shocks de precios de metales y estructura del empleo regional:\\
un modelo de factores específicos con margen formal--informal}
\author{}
\date{Julio 2026}

\begin{document}
\maketitle

\begin{abstract}
\noindent Este documento propone un \emph{marco teórico tentativo} —pensado como ejemplo de
estructura lógica— que explica cómo un shock de precios internacionales de metales se transmite
al mercado de trabajo de una región extractiva y reorganiza el empleo entre sectores y entre
formalidad e informalidad. La lógica es: primero se fijan supuestos mínimos; luego se describe
un mecanismo económico en un modelo de \emph{factores específicos} de tres sectores con un
margen de formalidad endógeno; de ese modelo se derivan proposiciones (estática comparativa);
y de cada proposición se obtiene una hipótesis \emph{contrastable} que se mapea a una prueba
empírica concreta. El objetivo no es un modelo definitivo, sino mostrar el andamiaje
``supuestos $\to$ mecanismo $\to$ proposiciones $\to$ hipótesis $\to$ contraste''.
\end{abstract}

\tableofcontents

\section{Qué hace un marco teórico y cómo está organizado este}
Un marco teórico cumple dos funciones: (i) \textbf{disciplinar} la relación entre variables
—decir por qué esperamos que $X$ afecte a $Y$ y con qué signo— y (ii) \textbf{generar
predicciones falsables} que el análisis empírico pueda confirmar o refutar. La estructura
lógica que sigue este documento es:
\[
\underbrace{\text{Supuestos}}_{\text{qué asumimos}}
\;\Rightarrow\;
\underbrace{\text{Mecanismo (modelo)}}_{\text{por qué se mueven las variables}}
\;\Rightarrow\;
\underbrace{\text{Proposiciones}}_{\text{qué implica el modelo}}
\;\Rightarrow\;
\underbrace{\text{Hipótesis}}_{\text{qué se contrasta}}.
\]
El lector puede reemplazar el contenido de cada bloque por el suyo; lo que importa es que la
cadena sea coherente y que cada hipótesis se derive de una proposición, no del aire.

\section{Hechos estilizados que el marco debe explicar}
El modelo se construye para ser consistente con cuatro regularidades documentadas:
\begin{enumerate}[label=(E\arabic*), leftmargin=2.2em, itemsep=1pt]
  \item La minería es \textbf{intensiva en capital y en el recurso}, con baja absorción directa
  de trabajo; sus efectos locales llegan sobre todo por \textbf{encadenamientos} hacia sectores
  no mineros \citep{aragon2013}.
  \item Los booms elevan el ingreso y reducen la pobreza local, pero con efectos que pueden ser
  \textbf{transitorios} (auge--caída) y espacialmente desiguales \citep{aragon2020,loayza2016}.
  \item La \textbf{informalidad} es el principal margen de ajuste del empleo ante shocks
  comerciales, absorbiendo trabajadores desplazados del sector formal \citep{dixcarneiro2017}.
  \item Existe un \textbf{canal fiscal} (canon minero) por el que el shock afecta el gasto
  público local, distinto del canal de mercado laboral \citep{loayza2016}.
\end{enumerate}

\section{Supuestos del modelo}

Considérese una región $r$ pequeña dentro de una economía abierta, poblada por $\bar{L}_r$
trabajadores. La región produce en tres sectores.

\begin{supuesto}[Estructura sectorial]\label{sup:sectores}
Hay tres sectores: minería $M$ (transable, precio internacional $p$), otros transables $T$
(precio internacional normalizado a 1) y no transables $N$ (precio local endógeno $q$, p.\ ej.\
comercio, servicios, construcción).
\end{supuesto}

\begin{supuesto}[Factores específicos]\label{sup:factores}
El sector $M$ usa un recurso fijo (el yacimiento) y capital, y emplea trabajo de forma
relativamente \emph{inelástica}. El trabajo es móvil \emph{entre sectores dentro} de la región,
pero \textbf{imperfectamente móvil entre regiones} (lo que justifica el análisis regional).
\end{supuesto}

\begin{supuesto}[Margen de formalidad]\label{sup:formalidad}
En los sectores $T$ y $N$ cada puesto puede ser formal o informal. La formalidad tiene un costo
$\tau>0$ (impuestos laborales, regulación) y un beneficio (acceso a capital/escala,
cumplimiento). Un puesto es formal si el excedente de formalizar es positivo; en el margen, la
proporción informal depende del salario y de $\tau$.
\end{supuesto}

\begin{supuesto}[Exposición y shock exógeno]\label{sup:shock}
La exposición de la región al metal es una participación \textbf{pre-determinada} $s_r\in[0,1]$
(peso de la minería en su economía de base). El precio $p$ se fija en el mercado mundial y es
\textbf{exógeno} a la región: $\bar{L}_r$ y $s_r$ no lo mueven.
\end{supuesto}

\begin{supuesto}[Canal fiscal]\label{sup:canon}
Una fracción de la renta minera se transfiere al gobierno regional (canon) $C_r=\phi\, s_r\, p$,
que financia gasto local (parte de él, demanda de trabajo formal).
\end{supuesto}

\section{El mecanismo: del precio mundial al empleo local}

\subsection{Ingreso regional y demanda de no transables}
La renta generada por la minería en la región es creciente en el precio y en la exposición:
\begin{equation}
Y_r \;=\; \underbrace{\pi_M(p)\, s_r}_{\text{renta minera}}
\;+\; \underbrace{w_r\,\bar{L}_r}_{\text{ingreso laboral}},
\qquad \frac{\partial \pi_M}{\partial p}>0.
\end{equation}
Como los transables tienen precio dado, el gasto adicional inducido por el boom se vuelca
desproporcionadamente sobre los \textbf{no transables} (cuyo precio sí puede subir). La demanda
de $N$ es creciente en el ingreso: $D_N = D_N(q, Y_r)$, con $\partial D_N/\partial Y_r>0$.

\subsection{Salario y precio local en equilibrio}
El empleo en $N$ satisface $w_r = q\cdot \mathrm{PMg}_L^N$. El vaciado del mercado de no
transables determina $q$; sustituyendo, el salario regional de equilibrio $w_r$ responde al
shock a través del ingreso. Formalmente, un aumento de $p$ eleva $Y_r$ en proporción a $s_r$,
lo que desplaza la demanda de $N$, sube $q$ y con ello $w_r$:
\begin{equation}
\frac{\partial w_r}{\partial p} \;=\; \underbrace{s_r}_{\text{exposición}}\cdot
\underbrace{\Psi}_{>0}\;>\;0,
\label{eq:dwdp}
\end{equation}
donde $\Psi>0$ recoge la fuerza del encadenamiento y la elasticidad de la oferta de $N$. La
ecuación~\eqref{eq:dwdp} es el corazón del mecanismo: \emph{el shock mundial se transmite al
salario local en proporción a la exposición regional}.

\subsection{La decisión de formalidad}
Sea $\theta_r$ la proporción de empleo informal en la región. Un firm-match se formaliza si su
excedente de formalidad $\Sigma$ supera cero. Ese excedente depende del salario (escala del
negocio), del costo $\tau$ y de la fase del ciclo:
\begin{equation}
\theta_r \;=\; \Theta\big(\,w_r,\ \tau,\ \text{fase}\,\big).
\end{equation}
El signo de $\partial \theta_r/\partial w_r$ es \textbf{ambiguo} y ese es el punto teórico
interesante:
\begin{itemize}[leftmargin=1.6em, itemsep=1pt]
  \item \textbf{Canal de formalización (signo $-$):} el boom expande la demanda; las firmas
  formales, con acceso a crédito y escala, crecen más rápido y absorben trabajadores en empleo
  formal $\Rightarrow$ $\theta_r$ baja.
  \item \textbf{Canal de entrada informal (signo $+$):} el boom atrae trabajadores y microu-
  nidades que entran \emph{informalmente} para capturar rentas rápido (autoempleo, minería
  artesanal) antes de que las firmas formales se ajusten $\Rightarrow$ $\theta_r$ sube.
\end{itemize}
Cuál domina es una \emph{cuestión empírica}: el modelo predice que el efecto neto sobre la
informalidad no está determinado a priori, y esa indeterminación es una predicción en sí misma.

\subsection{Asimetría auge--caída}
Introduzca fricciones realistas: costos de despido en el sector formal y rigidez a la baja de
salarios. Entonces la respuesta a $\Delta p>0$ y a $\Delta p<0$ no es simétrica:
\begin{itemize}[leftmargin=1.6em, itemsep=1pt]
  \item En el \textbf{auge}, el empleo formal se crea con fricción (contratación gradual) y
  parte del ajuste ocurre vía entrada informal.
  \item En la \textbf{caída}, el empleo formal se destruye y los trabajadores desplazados se
  refugian en el \textbf{autoempleo informal} como amortiguador \citep{dixcarneiro2017}.
\end{itemize}
Luego el aumento de $\theta_r$ en las caídas es mayor, en valor absoluto, que su variación en
los auges: la informalidad \textbf{amortigua asimétricamente}.

\section{Proposiciones (estática comparativa)}

\begin{prop}[Transmisión salarial proporcional a la exposición]\label{prop:w}
Bajo los Supuestos~\ref{sup:sectores}--\ref{sup:shock}, un aumento del precio del metal eleva
el salario local en proporción a la exposición de la región:
$\partial w_r/\partial p = s_r\,\Psi>0$. Regiones con $s_r=0$ no responden.
\end{prop}

\begin{prop}[Reasignación sectorial hacia no transables]\label{prop:reasig}
El empleo se reasigna desde $T$ hacia $N$ (y marginalmente hacia $M$), porque el boom eleva el
precio relativo de los no transables. La composición sectorial del empleo cambia aun cuando el
empleo total varíe poco.
\end{prop}

\begin{prop}[Efecto ambiguo sobre la informalidad]\label{prop:info}
El signo de $\partial \theta_r/\partial p$ depende de la fuerza relativa del canal de
formalización frente al de entrada informal; no está determinado por la teoría y debe
estimarse.
\end{prop}

\begin{prop}[Asimetría]\label{prop:asim}
Con costos de despido y rigidez salarial a la baja, la respuesta de la informalidad a una caída
de precios es mayor (en magnitud) que a un alza equivalente:
$\big|\partial \theta_r/\partial p\big|_{\Delta p<0} > \big|\partial \theta_r/\partial
p\big|_{\Delta p>0}$.
\end{prop}

\begin{prop}[Canal fiscal separable]\label{prop:canon}
Parte del efecto total sobre el empleo formal opera vía canon $C_r=\phi\,s_r\,p$ y no vía
mercado laboral privado. Controlar por $C_r$ aísla el canal de mercado laboral.
\end{prop}

\section{De las proposiciones a las hipótesis contrastables}

Cada hipótesis se deriva de una proposición y se empareja con una prueba empírica (que se
implementa con el diseño \textit{shift-share} descrito en el documento de estrategia de
identificación).

\begin{hipotesis}[Transmisión salarial]\label{h:w}
Un shock positivo de precios de metales incrementa el ingreso laboral de los hogares en las
regiones más expuestas. \emph{(De la Prop.~\ref{prop:w}; test: $\beta>0$ en la regresión del
log-ingreso sobre el índice Bartik $B_{rt}$.)}
\end{hipotesis}

\begin{hipotesis}[Reasignación sectorial]\label{h:reasig}
El shock aumenta la participación del empleo en no transables (comercio, servicios,
construcción) respecto a otros transables, en las regiones expuestas. \emph{(De la
Prop.~\ref{prop:reasig}; test: modelo multinomial de la celda sectorial.)}
\end{hipotesis}

\begin{hipotesis}[Informalidad, signo abierto]\label{h:info}
El efecto del shock sobre la probabilidad de empleo informal es distinto de cero, con signo a
determinar; se espera heterogeneidad entre gran minería (formalizadora) y zonas de minería
artesanal (informalizadora). \emph{(De la Prop.~\ref{prop:info}; test: $\beta$ sobre indicador
de informalidad, con interacciones por tipo de minería.)}
\end{hipotesis}

\begin{hipotesis}[Asimetría auge--caída]\label{h:asim}
La informalidad responde más a las caídas de precios que a los auges. \emph{(De la
Prop.~\ref{prop:asim}; test: $\beta^{+}\neq|\beta^{-}|$ separando la parte positiva y negativa
del shock.)}
\end{hipotesis}

\begin{hipotesis}[Canal fiscal]\label{h:canon}
Parte del efecto sobre el empleo formal desaparece al controlar por transferencias de canon,
evidenciando un canal fiscal distinto del laboral. \emph{(De la Prop.~\ref{prop:canon}; test:
cambio en $\beta$ al incluir el canon per cápita $C_{rt}$.)}
\end{hipotesis}

\section{La lógica causal, en un diagrama}
\begin{center}
\fbox{\parbox{0.94\textwidth}{\centering
$\Delta p^{\text{mundial}}_{\text{metal}} \times s_r^{\text{exposición base}}$
\;$\longrightarrow$\; $\uparrow$ renta e ingreso regional \;$\longrightarrow$\;
$\uparrow$ demanda de no transables \;$\longrightarrow$\; $\uparrow$ salario local\\[4pt]
$\longrightarrow$ \; \{ reasignación sectorial \;$+$\; margen formal/informal ambiguo \} \\[4pt]
con \; \textbf{asimetría} (fricciones de despido/rigidez) \; y \; \textbf{canal fiscal paralelo}
(canon)
}}
\end{center}

\section{Alcances y limitaciones del marco (honestidad teórica)}
\begin{itemize}[leftmargin=1.6em, itemsep=1pt]
  \item Es un modelo de \emph{estática comparativa} con fricciones añadidas de forma reducida;
  no resuelve una dinámica completa de búsqueda. Para la asimetría, una extensión natural es un
  modelo de búsqueda con destrucción de empleo formal.
  \item Trata la migración interregional como imperfecta pero exógena; endogeneizarla es una
  extensión relevante (y una amenaza empírica ya discutida en la estrategia de identificación).
  \item El margen de formalidad está modelado de forma reducida; podría microfundarse con una
  decisión firma-por-firma de cumplimiento y enforcement estocástico.
  \item No se modela explícitamente la heterogeneidad de calificación del trabajador, que
  \citet{benguria2024} muestran relevante; es una extensión directa (dos tipos de trabajo).
\end{itemize}
Estas limitaciones no invalidan las hipótesis: son la ``agenda'' de refinamiento teórico y
muestran dónde el marco puede crecer.

\begin{thebibliography}{9}
\setlength{\itemsep}{2pt}
\bibitem[Aragón y Rud(2013)]{aragon2013}
Aragón, F. M. y Rud, J. P. (2013). Natural Resources and Local Communities: Evidence from a
Peruvian Gold Mine. \textit{American Economic Journal: Economic Policy}, 5(2), 1--25.
\bibitem[Aragón y Rud(2020)]{aragon2020}
Aragón, F. M. y Rud, J. P. (2020). Fading Local Effects: Boom and Bust Evidence from a Peruvian
Gold Mine. \textit{Environment and Development Economics}, 25(2), 182--203.
\bibitem[Benguria et al.(2024)]{benguria2024}
Benguria, F., Saffie, F. y Zhang, S. (2024). Spatial Linkages and the Uneven Effects of a
Commodity Boom. \textit{Working Paper}.
\bibitem[Dix-Carneiro y Kovak(2017)]{dixcarneiro2017}
Dix-Carneiro, R. y Kovak, B. K. (2017). Trade Liberalization and Regional Dynamics.
\textit{American Economic Review}, 107(10), 2908--2946.
\bibitem[Loayza y Rigolini(2016)]{loayza2016}
Loayza, N. y Rigolini, J. (2016). The Local Impact of Mining on Poverty and Inequality:
Evidence from the Commodity Boom in Peru. \textit{World Development}, 84, 219--234.
\end{thebibliography}

\end{document}
